% \input{\pSections/"sec-staging.tex"}
\typeout{\pSections/"sec-staging.tex"}

\section{Staging}

%%    %%    %%    %%    %%    %%    %%    %%    %%    %%
\subsection{Images}

%%    %%    %%    %%    %%    %%    %%    %%    %%    %%
\subsection{Topics}
%
\begin{frame}\frametitle{\yl{Vector Tools} for Nonlinear Least Squares}
\begin{enumerate}
	\item \mfortran: \href{\urlMinpack}{minpack}, \href{https://people.math.sc.edu/Burkardt/f_src/nl2sol/nl2sol.html}{\texttt{NL2SOL}}, $\dots$
	\item \julia: \href{https://juliapackages.com/p/leastsquaresoptim}{LeastSquaresOptim.jl}
	\item \mathematica: \href{https://reference.wolfram.com/language/ref/FindMinimum.html}{\texttt{FindMinimum}}
	\item \octave: \href{https://octave.sourceforge.io/optim/function/leasqr.html}{\texttt{leasqr}}
	\item \python: \href{https://docs.scipy.org/doc/scipy/reference/generated/scipy.optimize.curve_fit.html}{SciPy}
	\item \href{https://petsc.org/release/}{\bl{PETSc}}: \href{https://petsc.org/release/manualpages/Tao/TAOBRGN/}{TAOBRGN}
\end{enumerate}
\end{frame}
%
\begin{frame}\frametitle{Nonlinear Least Squares References}
\begin{enumerate}
	\item \bl{LANL Briefing}: \href{https://mads.lanl.gov/presentations/Leif_LM_presentation_m.pdf}{Numerical Optimization using the
Levenberg-Marquardt Algorithm}\\[10pt]
	\item \bl{Duke}: \href{https://people.duke.edu/~hpgavin/ExperimentalSystems/lm.pdf}{The Levenberg-Marquardt algorithm}\\[10pt]
	\item \bl{\href{https://numerical.recipes/}{Numerical Recipes}}: \href{https://www.foo.be/docs-free/Numerical_Recipe_In_C/c15-5.pdf}{Nonlinear Models}\\[10pt]
	\item \bl{Wikipedia}: \href{https://en.wikipedia.org/wiki/Levenberg\%E2\%80\%93Marquardt_algorithm}{Levenberg–Marquardt algorithm}\\[10pt]
	\item \bl{Mathworld}: \href{https://mathworld.wolfram.com/Levenberg-MarquardtMethod.html}{Levenberg-Marquardt Method}
\end{enumerate}
\end{frame}

%%    %%    %%    %%    %%    %%    %%    %%    %%    %%
\subsection{Librarian Magic: Finding Sedov's Paper}

%\begin{frame}\frametitle{\href{https://en.wikipedia.org/wiki/Seshat}{Can't Find Seminal Work}}
%\center
%	\href{https://www.pinterest.com/pin/740982944913914922/}{
%	\begin{overpic}[ scale = 0.3 ]
%		{\pLocalGraphics help}
%		%\put(50, 50)	{\colorbox{white}{$a+b$}}
%	\end{overpic}}
%\end{frame}

%\begin{frame}\frametitle{\href{hhttps://www.britannica.com/topic/Seshat}{Seshat: Goddess of Writing, Wisdom, Knowledge}}
%\center
%	\href{https://www.pinterest.com/pin/740982944913914922/}{
%	\begin{overpic}[ scale = 0.225 ]
%		{\pLocalGraphics seshat}
%		%\put(50, 50)	{\colorbox{white}{$a+b$}}
%	\end{overpic}}
%\end{frame}

\begin{frame}\frametitle{Librarian Magic}
\small{
\begin{quotation}
They don't have enough information; neither did I. At the time of publication, it was not in English. \mg{Science Direct} doesn't offer this journal as far back as you've requested, under the name you cited. At the time, the journal was known as ``\bl{Prikladnaya matematika i mekhanika}'' with print \mg{ISSN 0032-8235}. 
\end{quotation}
}
\ 
\\[10pt]
\center
\href{https://translate.google.com/?sl=en&tl=ru&text=Prikladnaya\%20matematika\%20i\%20mekhanika\&op=translate}{\mg{\foreignlanguage{Russian}{Прикладная математика и механика}}}
\end{frame}
%
\begin{frame}\frametitle{Librarian Magic}
\small{
\begin{quotation}
That should help your UNM librarians, because according to \mg{WorldCat}, they will still have to request it from another library. There are \bl{only 4 in the US and Canada} who hold the title for the year you want... \mg{OCLC} charges per request initiated, so they're not putting incomplete information to get that request rejected, and have to submit (and pay for) another request. 
\end{quotation}
}
\end{frame}
%

\begin{frame}\frametitle{Librarian Magic\jumpLittle}
\center
	%\href{}{
	\begin{overpic}[ scale = 0.225 ]
		{\pLocalGraphics "JMM 1946 request"}
		%\put(50, 50)	{\colorbox{white}{$a+b$}}
	\end{overpic}
\end{frame}

%%    %%    %%    %%    %%    %%    %%    %%    %%    %%
\begin{frame}\frametitle{Librarian Magic}
\small{
\begin{quotation}
If you are ok receiving it in the original Russian, let them know. Also, academic libraries tend to close between Christmas and New Year so you may want to hold off on the ILL, lest it time out. 
\end{quotation}
}
\end{frame}

\endinput  %  ==  ==  ==  ==  ==  ==  ==  ==  ==
